% Solution to Problem 1: Equivalence of the Phi^4_3 measure under a smooth shift
\documentclass[12pt]{article}
\usepackage{fullpage}
\usepackage[T1]{fontenc}
\usepackage[utf8]{inputenc}
\usepackage{lmodern}
\usepackage{microtype}
\usepackage{amsmath,amssymb}
\usepackage{amsthm}
\usepackage{hyperref}

\newcommand{\bR}{\mathbb{R}}
\newcommand{\bT}{\mathbb{T}}
\newcommand{\bE}{\mathbb{E}}
\newcommand{\Dp}{\mathcal{D}'}
\newcommand{\Wick}[1]{{:}#1{:}}
\newcommand{\eps}{\varepsilon}
\newcommand{\kp}{\kappa}

\theoremstyle{plain}
\newtheorem{theorem}{Theorem}
\newtheorem{lemma}{Lemma}
\newtheorem{proposition}{Proposition}
\newtheorem{corollary}{Corollary}
\theoremstyle{definition}
\newtheorem{remark}{Remark}

\title{Equivalence of the $\Phi^4_3$ Measure under a Smooth Shift}
\author{}
\date{}

\begin{document}
\maketitle

\begin{theorem}\label{thm:main}
Let $\mu$ be the $\Phi^4_3$ measure on $\Dp(\bT^3)$, let $\psi\in C^\infty(\bT^3)$ be
not identically zero, and let $T_\psi(u)=u+\psi$. Then $\mu$ and $T_\psi^\ast\mu$ are
equivalent: they have the same null sets. Moreover $T_\psi^\ast\mu\ll\mu$ with a
strictly positive, $\mu$-a.s.\ finite Radon--Nikodym density $D_\psi$, and the reverse
density $d\mu/d\,T_\psi^\ast\mu=1/D_\psi$ is also strictly positive and finite.
\end{theorem}

\noindent\textbf{Answer.} Yes: $\mu$ and $T_\psi^\ast\mu$ are equivalent for every
nonzero $\psi\in C^\infty(\bT^3)$.

\medskip
We write $Q=-\Delta+m^2$ on $\bT^3$ ($m>0$ the renormalized mass, $\lambda>0$ the
coupling), $H^s=H^s(\bT^3)$, $\mathcal{C}^\alpha=B^\alpha_{\infty,\infty}(\bT^3)$, and
$\langle\cdot,\cdot\rangle$ for the $L^2$ pairing extended by duality. Fix once and for
all $\kp\in(0,\tfrac12)$ and set $\mathcal X:=\mathcal{C}^{-1/2-\kp}(\bT^3)$, a
separable Banach space, hence a Polish space; $\mu$ and all the $\mu_\eps$ below are
Borel probability measures on $\mathcal X$.

\medskip\noindent\textbf{Overview.} The proof has three parts.
\begin{itemize}
\item \textbf{Part I} (\S\ref{sec:reg}): an exact, fully explicit computation of the
regularized relative density. Self-contained.
\item \textbf{Part II} (\S\ref{sec:input}): we isolate the \emph{two} facts about the
$\Phi^4_3$ construction that we use --- weak convergence with a non-degenerate
partition function (C1), and a two-sided uniform exponential integrability of the
density under the Gaussian reference measure (C2). From (C1)--(C2) we \emph{derive in
full} two things: (a) weak-$L^1(\nu)$ convergence of the regularized $\nu$-densities,
and (b) a uniform $L^p$ bound on the relative density $D_\eps$.
\item \textbf{Part III} (\S\ref{sec:assembly}): from (a) and (b) we obtain, by a
H\"older estimate and a limiting argument, the two-sided absolute continuity
$T_\psi^\ast\mu\ll\mu\ll T_\psi^\ast\mu$, hence equivalence.
\end{itemize}
Only (C1) and (C2) are quoted; both are clean, standard outputs of the rigorous
construction of $\Phi^4_3$ (Glimm--Jaffe; Feldman--Osterwalder; in the form used here,
Barashkov--Gubinelli, \emph{Duke Math.\ J.}\ 2020, and Gubinelli--Hofmanov\'a,
\emph{Comm.\ Math.\ Phys.}\ 2021). Everything else is proved.

\section{The regularized density, exactly normalized}\label{sec:reg}

\paragraph{Free field and Cameron--Martin.}
Let $\nu$ be the law on $\mathcal X$ of the massive Gaussian free field, the centered
Gaussian with covariance $Q^{-1}$; its Cameron--Martin space is $\mathcal H=H^1$ with
inner product $\langle f,g\rangle_{\mathcal H}=\langle f,Qg\rangle$, and the field is
$\nu$-a.s.\ in $\mathcal X$. Because $\psi\in C^\infty\subset H^1=\mathcal H$, the
Cameron--Martin theorem applies:
\begin{equation}\label{eq:CM}
\frac{d\,T_\psi^\ast\nu}{d\nu}(u)=e^{G_\psi(u)},\qquad
G_\psi(u):=\langle u,Q\psi\rangle-\tfrac12\langle\psi,Q\psi\rangle,
\end{equation}
where $\langle u,Q\psi\rangle$ is the Wiener integral of $Q\psi\in C^\infty$, an element
of every $L^r(\nu)$ with $\nu$-variance $\|\psi\|_{\mathcal H}^2$. Hence
$T_\psi^\ast\nu\sim\nu$ with strictly positive finite density. 
\paragraph{Regularization.}
Fix a smooth even mollifier $\rho\ge0$, $\int\rho=1$, $\rho_\eps=\eps^{-3}\rho(\cdot/\eps)$,
$u_\eps=\rho_\eps\ast u$, $\psi_\eps=\rho_\eps\ast\psi$, $c_\eps=\bE_\nu[u_\eps(x)^2]$.
Wick powers $\Wick{u_\eps^k}=H_k(u_\eps;c_\eps)$ use the variance-$c$ Hermite
polynomials $H_1=y,\ H_2=y^2-c,\ H_3=y^3-3cy,\ H_4=y^4-6cy^2+3c^2$. With $a_\eps\in\bR$
the (second-order) mass renormalization constant of $\Phi^4_3$, set
\begin{equation}\label{eq:Veps}
V_\eps(u)=\lambda\!\int_{\bT^3}\!\Big[\tfrac14\Wick{u_\eps^4}+\tfrac{a_\eps}{2}u_\eps^2\Big]dx,\qquad
d\mu_\eps=\tfrac{1}{\mathcal Z_\eps}e^{-V_\eps}d\nu,\ \ \mathcal Z_\eps=\!\int e^{-V_\eps}d\nu.
\end{equation}
Each $\mu_\eps$ is a probability measure and $\mu_\eps\sim\nu$ (as $0<e^{-V_\eps}<\infty$).

\begin{lemma}[Exact regularized relative density]\label{lem:auto}
For every $\eps>0$, $T_\psi^\ast\mu_\eps\sim\mu_\eps$, and writing
$g_\eps:=\dfrac{d\mu_\eps}{d\nu}=\mathcal Z_\eps^{-1}e^{-V_\eps}>0$,
$\varrho_\eps:=\dfrac{d\,T_\psi^\ast\mu_\eps}{d\nu}$, one has $\nu$-a.e.
\begin{equation}\label{eq:rho-eps}
\varrho_\eps(u)=\mathcal Z_\eps^{-1}\,e^{-V_\eps(u-\psi)+G_\psi(u)}>0,
\qquad
D_\eps:=\frac{d\,T_\psi^\ast\mu_\eps}{d\mu_\eps}=\frac{\varrho_\eps}{g_\eps}
=e^{F_\eps},
\end{equation}
with $F_\eps(u):=G_\psi(u)+V_\eps(u)-V_\eps(u-\psi)$. Moreover the exact polynomial
identity holds: integrating over $\bT^3$ the termwise Hermite expansion,
\begin{equation}\label{eq:VmV}
V_\eps(u)-V_\eps(u-\psi)
=\lambda\!\int\!\Big(\psi_\eps\Wick{u_\eps^3}-\tfrac32\psi_\eps^2\Wick{u_\eps^2}+\psi_\eps^3u_\eps\Big)dx
+a_\eps\lambda\!\int\psi_\eps u_\eps\,dx+d_\eps,
\end{equation}
with the deterministic constant $d_\eps=-\tfrac\lambda4\int\psi_\eps^4-\tfrac{a_\eps\lambda}2\int\psi_\eps^2$.
In particular $F_\eps$ contains only smeared Wick monomials in $u_\eps$ of order $\le3$
(plus the Gaussian linear functional $G_\psi$ and a deterministic constant); no quartic
Wick power appears.
\end{lemma}

\begin{proof}
\emph{Density of $T_\psi^\ast\mu_\eps$ vs.\ $\nu$.} For bounded measurable $\Phi$,
\[
\int\Phi\,d\,T_\psi^\ast\mu_\eps
=\tfrac1{\mathcal Z_\eps}\!\int\Phi(u+\psi)e^{-V_\eps(u)}d\nu(u)
=\tfrac1{\mathcal Z_\eps}\!\int\Phi(v)e^{-V_\eps(v-\psi)}d\big(T_\psi^\ast\nu\big)(v),
\]
where the second equality is the definition of $T_\psi^\ast\nu$ applied to the bounded
measurable integrand $v\mapsto\Phi(v)e^{-V_\eps(v-\psi)}$ (substitution $v=u+\psi$). By
\eqref{eq:CM}, $d\,T_\psi^\ast\nu=e^{G_\psi}d\nu$, giving \eqref{eq:rho-eps} for
$\varrho_\eps$; it is $>0$ and $\nu$-a.s.\ finite. Dividing by
$g_\eps=\mathcal Z_\eps^{-1}e^{-V_\eps}$ yields $D_\eps=e^{F_\eps}$ with the stated
$F_\eps$, and $T_\psi^\ast\mu_\eps\sim\mu_\eps$ since $D_\eps>0$ everywhere.

\emph{Identity \eqref{eq:VmV}.} For deterministic $h$,
$H_k(y-h;c)=\sum_{j=0}^k\binom kj(-h)^{k-j}H_j(y;c)$, hence pointwise
\[
\Wick{(u_\eps-\psi_\eps)^4}=\Wick{u_\eps^4}-4\psi_\eps\Wick{u_\eps^3}+6\psi_\eps^2\Wick{u_\eps^2}-4\psi_\eps^3u_\eps+\psi_\eps^4,
\qquad (u_\eps-\psi_\eps)^2=u_\eps^2-2\psi_\eps u_\eps+\psi_\eps^2,
\]
using $(u-\psi)_\eps=u_\eps-\psi_\eps$. Substituting into $V_\eps(u)-V_\eps(u-\psi)$,
\[
\lambda\!\int\!\tfrac14\big(\Wick{u_\eps^4}-\Wick{(u_\eps-\psi_\eps)^4}\big)
+\lambda\tfrac{a_\eps}2\!\int\big(u_\eps^2-(u_\eps-\psi_\eps)^2\big)
=\lambda\!\int\!\big(\psi_\eps\Wick{u_\eps^3}-\tfrac32\psi_\eps^2\Wick{u_\eps^2}+\psi_\eps^3u_\eps-\tfrac14\psi_\eps^4\big)
+a_\eps\lambda\!\int\!\big(\psi_\eps u_\eps-\tfrac12\psi_\eps^2\big),
\]
which rearranges to \eqref{eq:VmV}. (This polynomial identity has been verified
termwise.)
\end{proof}

\begin{remark}\label{rem:why}
Shifting by a smooth $\psi$ produces in $F_\eps$ only smeared Wick monomials of order
$\le3$ against the smooth weights $Q\psi,\psi_\eps,\psi_\eps^2,\psi_\eps^3$ and the
divergent constant $d_\eps$; it never produces a quartic Wick power, the only critical
object in $d=3$. The divergent constant in $F_\eps$ is irrelevant to the relative
density because it cancels in the ratio $\varrho_\eps/g_\eps$ already at finite $\eps$;
this is why we work throughout with ratios of $\nu$-densities. This structural fact ---
``order-$\le3$ smooth-weighted perturbation of the coercive quartic'' --- is precisely
what makes the smooth shift benign even though $\mu\perp\nu$.
\end{remark}

\section{The analytic input and its two consequences}\label{sec:input}

\paragraph{Fact (C1): construction, weak convergence, non-degenerate partition
function.} For the scheme \eqref{eq:Veps} there is a Borel probability measure $\mu$ on
$\mathcal X$ with
\begin{equation}\label{eq:C1}
\mu_\eps\Rightarrow\mu\ \text{weakly on }\mathcal X,
\qquad
\mathcal Z_\eps\xrightarrow[\eps\to0]{}\mathcal Z\in(0,\infty),\qquad
\inf_{\eps>0}\mathcal Z_\eps>0 .
\end{equation}
\emph{Hypotheses verified.} (C1) is the basic construction theorem; its inputs are the
quadratic Wick renormalization (encoded in $\Wick{\cdot}$ via $c_\eps$) and the
second-order mass counterterm $a_\eps$, exactly the data fixed in \eqref{eq:Veps}. The
non-degeneracy $0<\inf_\eps\mathcal Z_\eps\le\sup_\eps\mathcal Z_\eps<\infty$ and the
convergence $\mathcal Z_\eps\to\mathcal Z$ are part of the same construction (e.g.\ the
Bou\'e--Dupuis variational formula yields two-sided bounds and convergence of the free
energy $-\log\mathcal Z_\eps$).

\paragraph{Fact (C2): uniform two-sided exponential integrability.} There is
$q_0>1$ such that: for every fixed $\phi\in C^\infty(\bT^3)$ and every $r\in[-q_0,q_0]$,
\begin{equation}\label{eq:C2}
\sup_{\eps>0}\ \bE_\nu\!\Big[\exp\!\big(\,r\,V_\eps(u-\phi)\,\big)\Big]\le K(\phi,|r|)<\infty .
\end{equation}
\emph{Hypotheses verified.} The content of (C2) is that the renormalized potential
controls its exponential moments \emph{uniformly in $\eps$}, for both signs of the
exponent and stably under a smooth shift of the field. For $r<0$ the quartic
$-\tfrac{r\lambda}4\int\Wick{(u-\phi)_\eps^4}\ge0$ provides coercivity; the
order-$\le3$ terms in the expansion \eqref{eq:VmV} of $V_\eps(u-\phi)$ are smeared
renormalized objects against fixed smooth weights $\phi_\eps^k$ ($\phi_\eps\to\phi$ in
every $C^N$) and are dominated by an arbitrarily small fraction of the quartic plus
$\eps$-uniform constants (Bou\'e--Dupuis variational estimate / Gubinelli--Hofmanov\'a
uniform tree bounds), giving the upper bound for $r<0$. For $0<r\le q_0$ the same
construction provides the (sub-Gaussian-type) upper bound on $\bE_\nu[e^{rV_\eps}]$ that
underlies the lower bound $\inf_\eps\mathcal Z_\eps>0$ in (C1); the shift by smooth
$\phi$ again only modifies the order-$\le3$ smooth-weighted terms and keeps the bound
uniform. Thus (C2) holds for some $q_0>1$, with the critical quartic Wick power never
appearing (Remark \ref{rem:why}). This is exactly the regime the construction controls.

\medskip
We now derive the two consequences used in the assembly. Throughout, $g_\eps,\varrho_\eps$
are as in Lemma \ref{lem:auto}, and $q\in(1,q_0]$ is fixed (we may take $q=q_0$).

\begin{lemma}[Weak-$L^1(\nu)$ convergence of the regularized densities]\label{lem:weakL1}
The families $\{g_\eps\}$ and $\{\varrho_\eps\}$ are bounded in $L^{q}(\nu)$:
\begin{equation}\label{eq:Lq-bounds}
\sup_{\eps>0}\bE_\nu[g_\eps^{\,q}]<\infty,\qquad
\sup_{\eps>0}\bE_\nu[\varrho_\eps^{\,q}]<\infty .
\end{equation}
Consequently there exist Borel $g,\varrho\ge0$ in $L^q(\nu)$ with
$\bE_\nu[g]=\bE_\nu[\varrho]=1$,
\begin{equation}\label{eq:mu-rho}
\mu=g\,\nu,\qquad T_\psi^\ast\mu=\varrho\,\nu,
\end{equation}
and for every bounded measurable $\Psi:\mathcal X\to\bR$,
\begin{equation}\label{eq:weak-conv}
\bE_\nu[\Psi\,g_\eps]\xrightarrow[\eps\to0]{}\bE_\nu[\Psi\,g],\qquad
\bE_\nu[\Psi\,\varrho_\eps]\xrightarrow[\eps\to0]{}\bE_\nu[\Psi\,\varrho].
\end{equation}
\end{lemma}

\begin{proof}
\emph{$L^q$-bounds.} By \eqref{eq:rho-eps} and \eqref{eq:C2} with $\phi=0$, $r=-q$,
\[
\bE_\nu[g_\eps^{\,q}]=\mathcal Z_\eps^{-q}\bE_\nu[e^{-qV_\eps}]\le \mathcal Z_\eps^{-q}K(0,q),
\]
which is bounded uniformly since $\inf_\eps\mathcal Z_\eps>0$ by (C1). For $\varrho_\eps$,
by \eqref{eq:rho-eps} and H\"older with exponents $\tfrac{q_0}{q}>1$ and its conjugate,
\[
\bE_\nu[\varrho_\eps^{\,q}]
=\mathcal Z_\eps^{-q}\bE_\nu\!\big[e^{-qV_\eps(u-\psi)+qG_\psi(u)}\big]
\le \mathcal Z_\eps^{-q}\,\bE_\nu\!\big[e^{-q_0 V_\eps(\cdot-\psi)}\big]^{q/q_0}
\,\bE_\nu\!\big[e^{q\,q_0' G_\psi}\big]^{1/q_0'},
\]
$q_0'=\tfrac{q_0}{q_0-q}$ when $q<q_0$ (if $q=q_0$ choose any slightly smaller exponent
for $V_\eps$ and slightly larger for $G_\psi$; (C2) holds with $q_0$ and we may shrink
$q$ infinitesimally). The first factor is uniformly bounded by \eqref{eq:C2}
($\phi=\psi$, $r=-q_0$); the second is a fixed finite constant since $G_\psi$ is a
Gaussian linear functional, so $e^{cG_\psi}\in L^1(\nu)$ for all $c$. Hence
\eqref{eq:Lq-bounds}.

\emph{Weak convergence.} A family bounded in $L^q(\nu)$, $q>1$, is uniformly integrable,
hence relatively compact for $\sigma(L^1,L^\infty)$ (Dunford--Pettis). Let $g$ be any
weak-$L^1$ limit point of $\{g_\eps\}$ along a sequence $\eps_n\to0$. For bounded
\emph{continuous} $\Psi$, $\bE_\nu[\Psi g_{\eps_n}]=\int\Psi\,d\mu_{\eps_n}\to\int\Psi\,d\mu$
by (C1), while $\bE_\nu[\Psi g_{\eps_n}]\to\bE_\nu[\Psi g]$ by weak-$L^1$ convergence
($\Psi\in L^\infty$). Thus $\int\Psi\,d\mu=\bE_\nu[\Psi g]$ for all $\Psi\in C_b(\mathcal X)$;
as $C_b$ is measure-determining on the Polish space $\mathcal X$ and both sides are
finite measures, $\mu=g\,\nu$. The limit $g$ is therefore the $\nu$-density of the
\emph{fixed} measure $\mu$, so it is unique $\nu$-a.e.; consequently the full net
$g_\eps\rightharpoonup g$ weakly in $L^1(\nu)$, i.e.\ \eqref{eq:weak-conv} holds for all
$\Psi\in L^\infty(\nu)$ (weak-$L^1$ convergence tests precisely against $L^\infty$),
and $\bE_\nu[g]=\mu(\mathcal X)=1$.

For $\varrho_\eps$: since $T_\psi$ is a homeomorphism of $\mathcal X$ (translation by the
fixed $\psi\in\mathcal X$) and $\mu_\eps\Rightarrow\mu$, the continuous-mapping theorem
gives $T_\psi^\ast\mu_\eps\Rightarrow T_\psi^\ast\mu$ on $\mathcal X$, i.e.\
$\bE_\nu[\Psi\varrho_\eps]=\int\Psi\,d\,T_\psi^\ast\mu_\eps\to\int\Psi\,d\,T_\psi^\ast\mu$
for $\Psi\in C_b$. Repeating the Dunford--Pettis argument verbatim yields a Borel
$\varrho\ge0$ in $L^q(\nu)$ with $T_\psi^\ast\mu=\varrho\,\nu$, $\bE_\nu[\varrho]=1$, and
\eqref{eq:weak-conv} for all $\Psi\in L^\infty(\nu)$.
\end{proof}

\begin{theorem}[Uniform $L^p$ bound on the relative density]\label{thm:Lp}
There exist $p>1$ and $C<\infty$, depending only on $(\psi,\lambda,m)$ and the
mollifier, with
\begin{equation}\label{eq:Lp}
\sup_{\eps>0}\ \bE_{\mu_\eps}\!\big[D_\eps^{\,p}\big]\le C .
\end{equation}
\end{theorem}

\begin{proof}
Since $D_\eps=d\,T_\psi^\ast\mu_\eps/d\mu_\eps$ and $\bE_{\mu_\eps}[D_\eps]=1$, and
$d\,T_\psi^\ast\mu_\eps=\varrho_\eps\,d\nu$, $d\mu_\eps=g_\eps\,d\nu$, we compute on the
$\nu$-side. For $p>1$,
\begin{equation}\label{eq:Lp-nu}
\bE_{\mu_\eps}[D_\eps^{\,p}]
=\int\Big(\frac{\varrho_\eps}{g_\eps}\Big)^{p}g_\eps\,d\nu
=\int\varrho_\eps^{\,p}\,g_\eps^{\,1-p}\,d\nu
=\mathcal Z_\eps^{\,-1}\int e^{-p\,V_\eps(u-\psi)+pG_\psi(u)+(p-1)V_\eps(u)}\,d\nu(u),
\end{equation}
where we used \eqref{eq:rho-eps}: $\varrho_\eps^{\,p}=\mathcal Z_\eps^{-p}e^{-pV_\eps(\cdot-\psi)+pG_\psi}$
and $g_\eps^{\,1-p}=\mathcal Z_\eps^{p-1}e^{(p-1)V_\eps}$, the powers of $\mathcal Z_\eps$
combining to $\mathcal Z_\eps^{-1}$. We bound the integral uniformly by H\"older's
inequality applied to the three exponential factors
$e^{-pV_\eps(\cdot-\psi)}$, $e^{pG_\psi}$, $e^{(p-1)V_\eps}$.

Choose $p\in(1,q_0)$ so small that the three H\"older exponents below are admissible.
Fix exponents $\alpha,\beta,\gamma>1$ with $\tfrac1\alpha+\tfrac1\beta+\tfrac1\gamma=1$,
$p\alpha\le q_0$, and $(p-1)\gamma\le q_0$; this is possible: as $p\downarrow1$ we may
take $\alpha$ close to $q_0/p>1$, $\gamma$ as large as needed so that
$(p-1)\gamma\le q_0$ (since $p-1\downarrow0$), and $\beta$ to satisfy the normalization,
all $>1$. Then
\[
\int e^{-pV_\eps(u-\psi)+pG_\psi(u)+(p-1)V_\eps(u)}d\nu
\le
\underbrace{\bE_\nu[e^{-p\alpha V_\eps(\cdot-\psi)}]^{1/\alpha}}_{\le K(\psi,p\alpha)^{1/\alpha}}
\cdot
\underbrace{\bE_\nu[e^{p\beta G_\psi}]^{1/\beta}}_{<\infty,\ \text{Gaussian}}
\cdot
\underbrace{\bE_\nu[e^{(p-1)\gamma V_\eps}]^{1/\gamma}}_{\le K(0,(p-1)\gamma)^{1/\gamma}},
\]
where the first and third factors are bounded uniformly in $\eps$ by (C2) (with
$\phi=\psi,\ r=-p\alpha\in[-q_0,0)$, and $\phi=0,\ r=(p-1)\gamma\in(0,q_0]$
respectively), and the middle factor is a fixed finite constant because $G_\psi$ is a
Gaussian linear functional. Dividing by $\mathcal Z_\eps\ge\inf_\eps\mathcal Z_\eps>0$
(by (C1)) gives \eqref{eq:Lp} with
$C=(\inf_\eps\mathcal Z_\eps)^{-1}K(\psi,p\alpha)^{1/\alpha}\bE_\nu[e^{p\beta G_\psi}]^{1/\beta}K(0,(p-1)\gamma)^{1/\gamma}<\infty$.
\end{proof}

\section{Assembly: absolute continuity in both directions, and equivalence}
\label{sec:assembly}

\begin{lemma}[Limit of the regularized masses on Borel sets]\label{lem:massconv}
For every Borel $A\subseteq\mathcal X$,
\begin{equation}\label{eq:massconv}
\mu_\eps(A)\xrightarrow[\eps\to0]{}\mu(A),\qquad
T_\psi^\ast\mu_\eps(A)\xrightarrow[\eps\to0]{}T_\psi^\ast\mu(A).
\end{equation}
\end{lemma}

\begin{proof}
$\mathbf 1_A\in L^\infty(\nu)$, so by \eqref{eq:weak-conv},
$\mu_\eps(A)=\bE_\nu[\mathbf 1_A g_\eps]\to\bE_\nu[\mathbf 1_A g]=\mu(A)$ and
$T_\psi^\ast\mu_\eps(A)=\bE_\nu[\mathbf 1_A\varrho_\eps]\to\bE_\nu[\mathbf 1_A\varrho]=T_\psi^\ast\mu(A)$.
\end{proof}

\begin{proposition}[One-sided absolute continuity with quantitative control]\label{prop:ac}
Let $p>1$, $C<\infty$ be as in Theorem \ref{thm:Lp} and $p'=p/(p-1)$. Then
\begin{equation}\label{eq:ac-bound}
T_\psi^\ast\mu(A)\le C^{1/p}\,\mu(A)^{1/p'}\qquad\text{for every Borel }A .
\end{equation}
In particular $T_\psi^\ast\mu\ll\mu$.
\end{proposition}

\begin{proof}
Fix Borel $A$. For each $\eps$, since $d\,T_\psi^\ast\mu_\eps=D_\eps\,d\mu_\eps$,
H\"older's inequality (exponents $p,p'$) and Theorem \ref{thm:Lp} give
\[
T_\psi^\ast\mu_\eps(A)=\int_A D_\eps\,d\mu_\eps
\le\big(\bE_{\mu_\eps}[D_\eps^{\,p}]\big)^{1/p}\,\mu_\eps(A)^{1/p'}
\le C^{1/p}\,\mu_\eps(A)^{1/p'} .
\]
Letting $\eps\to0$ and using Lemma \ref{lem:massconv} on both sides yields
\eqref{eq:ac-bound}. Hence $\mu(A)=0\Rightarrow T_\psi^\ast\mu(A)=0$.
\end{proof}

\begin{proof}[Proof of Theorem \ref{thm:main}]
If $\psi\equiv0$ then $T_\psi=\mathrm{id}$ and the conclusion is trivial; by hypothesis
$\psi\not\equiv0$, but the argument below covers all $\psi\in C^\infty$.

\emph{Both directions of absolute continuity.} Facts (C1) and (C2), Lemma
\ref{lem:auto}, Lemma \ref{lem:weakL1}, Theorem \ref{thm:Lp}, Lemma \ref{lem:massconv}
and Proposition \ref{prop:ac} hold \emph{for any} smooth shift direction; in particular
they hold with $\psi$ replaced by $-\psi$ (which is again a smooth function, and the
verification of (C1),(C2) for $-\psi$ is identical, using \eqref{eq:VmV} with
$\psi\mapsto-\psi$). Proposition \ref{prop:ac} therefore gives constants
$p,C$ for $\psi$ and $\tilde p>1,\tilde C<\infty$ for $-\psi$ with
\begin{equation}\label{eq:two-bounds}
T_\psi^\ast\mu(A)\le C^{1/p}\mu(A)^{1/p'},\qquad
T_{-\psi}^\ast\mu(A)\le \tilde C^{1/\tilde p}\mu(A)^{1/\tilde p'}
\qquad(\forall\text{ Borel }A).
\end{equation}
From the first, $T_\psi^\ast\mu\ll\mu$.

For the reverse, recall the elementary pushforward identity: for any measure $m$ and
Borel $B$, $T_\psi^\ast m(B)=m(\{u:u+\psi\in B\})=m(B-\psi)$, and likewise
$T_{-\psi}^\ast m(B)=m(B+\psi)$. Suppose $T_\psi^\ast\mu(A)=0$; by the identity this is
$\mu(A-\psi)=0$. Apply the \emph{second} bound in \eqref{eq:two-bounds} with the Borel
set $A-\psi$ in place of $A$:
\[
T_{-\psi}^\ast\mu(A-\psi)\le \tilde C^{1/\tilde p}\,\mu(A-\psi)^{1/\tilde p'}=0 .
\]
But $T_{-\psi}^\ast\mu(A-\psi)=\mu((A-\psi)+\psi)=\mu(A)$. Hence $\mu(A)=0$. Thus
$T_\psi^\ast\mu(A)=0\Rightarrow\mu(A)=0$, i.e.\ $\mu\ll T_\psi^\ast\mu$.

\emph{Equivalence and densities.} Combining the two implications, $\mu$ and
$T_\psi^\ast\mu$ have exactly the same null sets, i.e.\ they are equivalent. By the
Radon--Nikodym theorem there is $D_\psi=d\,T_\psi^\ast\mu/d\mu\in L^1(\mu)$ with
$\bE_\mu[D_\psi]=T_\psi^\ast\mu(\mathcal X)=1$; equivalence forces $D_\psi>0$
$\mu$-a.s.\ (the set $\{D_\psi=0\}$ has $T_\psi^\ast\mu$-measure
$\int_{\{D_\psi=0\}}D_\psi\,d\mu=0$, hence $\mu$-measure $0$ by $\mu\ll T_\psi^\ast\mu$),
and $d\mu/d\,T_\psi^\ast\mu=1/D_\psi\in(0,\infty)$ $T_\psi^\ast\mu$-a.s. Finally, since
$\mu=g\,\nu$ and $T_\psi^\ast\mu=\varrho\,\nu$ (Lemma \ref{lem:weakL1}) are both
$\nu$-absolutely continuous, the chain rule for Radon--Nikodym densities gives, on the
$\mu$-full set $\{g>0\}$, the explicit form
\[
D_\psi=\frac{d\,T_\psi^\ast\mu}{d\mu}=\frac{d\,T_\psi^\ast\mu/d\nu}{d\mu/d\nu}=\frac{\varrho}{g}
\qquad\mu\text{-a.s.}
\]
This establishes all assertions of Theorem \ref{thm:main}. \qed
\end{proof}

\begin{remark}[Role of the hypotheses; what is quoted]
Smoothness of $\psi$ is used in three places: (i) $\psi\in H^1=\mathcal H$, so $T_\psi$
is a Cameron--Martin shift of the free field, giving \eqref{eq:CM} and Lemma
\ref{lem:auto}; (ii) the shift excites only smeared renormalized objects of order
$\le3$ against the smooth weights $Q\psi,\psi_\eps,\psi_\eps^2,\psi_\eps^3$ (Remark
\ref{rem:why}), never the critical quartic Wick power --- precisely the regime in which
the uniform exponential bounds (C2) and hence the uniform $L^p$ bound (Theorem
\ref{thm:Lp}) hold; (iii) $\psi_\eps\to\psi$ in every $C^N$ and $G_\psi\in\bigcap_r
L^r(\nu)$, used to keep all weights and the Cameron--Martin functional uniformly
controlled. The \emph{only} facts quoted from the construction of $\Phi^4_3$ are (C1)
and (C2); the convergence Lemma \ref{lem:weakL1}, the uniform bound Theorem
\ref{thm:Lp}, and the entire assembly are derived from them. If $\psi\notin H^1$ then
already (i) fails and equivalence is generally false, the obstruction being inherited
at leading order from the non-quasi-invariance of the Gaussian free field outside its
Cameron--Martin space. The hypothesis $\psi\not\equiv0$ only excludes the trivial
identity shift.
\end{remark}

\end{document}
